%======================================================================
\section{The Pauli configuration-lamp group}\label{sec:group}
%======================================================================

This section builds the group $\GP$ behind Theorems~A and~B and proves the facts about it that the rest of the paper uses: $J$ is a
central involution with $J\ne1$, the six elements $p,q,t,d,x,z$ generate $\GP$ (\cref{prop:J-central}), and $\GP$ is amenable
(\cref{prop:GP-amenable}). Its extension $\GPp$ by a ray swap has the same properties (\cref{lem:GPp}); the game of \cref{sec:game} is
built on it. The construction has three layers, drawn in \cref{fig:lamps}. The \emph{points} are those of the three rays of Houghton's
group $H_3$, which moves them by permutations that are eventually translations. The \emph{sites} are the finite sets of points, written
as binary vectors and called configurations; an affine group $A$ moves sites by permuting points, toggling a point, or adding one point
into another. Each site carries a \emph{Pauli pair}, a copy of the dihedral group $D_4\cong\langle\sigma_X,\sigma_Z\rangle$, and the
centres of all the pairs are identified to one element $J$. On $N$ points there are $2^N$ sites and hence $2^N$ Pauli pairs, any two of
which commute, while the generators $t$ and $d$ that change sites involve only one and two points. \Cref{sec:pair} shows that the commutation of any two
of these pairs follows from a fixed finite set of relations at a cost quasi-polynomial in $N$, and \cref{sec:sector} turns this into the
entanglement lower bound. That cost is measured by the Dehn function of a presentation of $H_3$, which we fix in
\cref{sec:group-houghton} together with the bound proved in \cref{app:dehn}. The host group $A$ and the generators $p,q,t,d$ are those
of~\cite[Section~3]{DM-lamps}, where the lamps are copies of $S_3$.

\subsection{Houghton's group}\label{sec:group-houghton}

\begin{definition}[Houghton's group]\label{def:houghton}
Let $Y=\{1,2,3\}\times\{1,2,\dots\}$; the point $(i,j)$ lies on \emph{ray}~$i$ at \emph{height}~$j$. \emph{Houghton's group}
$H=H_3$~\cite{Hou78} is the group of permutations $g$ of $Y$ that are eventually translations on each ray: there are integers
$m_1(g),m_2(g),m_3(g)$ with $g(i,j)=(i,j+m_i(g))$ for all large $j$. Let $p\in H$ be the permutation
\[
p(1,j)=(1,j+1),\qquad p(2,1)=(1,1),\qquad p(2,j)=(2,j-1)\ \ (j\ge2),\qquad p(3,j)=(3,j),
\]
let $q$ be defined in the same way with rays $2$ and $3$ exchanged, and put $s=[p,q]=pqp^{-1}q^{-1}$.
\end{definition}

The finitary permutations of $Y$ form a subgroup $\FSym(Y)\le H$. The map $g\mapsto(m_1(g),m_2(g),m_3(g))$ is a homomorphism to $\Z^3$,
and its kernel is $\FSym(Y)$, because an element with all $m_i(g)=0$ fixes every point of large height. Moreover
$m_1(g)+m_2(g)+m_3(g)=0$: if $g$ translates every point of height above $K$ on ray $i$ by $m_i(g)$, and $B$ is the set of the $3K$ points
of height at most $K$, then $g(B)$ is the complement of $g(Y\setminus B)$, which has $\sum_i(K+m_i(g))$ points. As $m(p)=(1,-1,0)$ and
$m(q)=(1,0,-1)$, recording $-m_2$ and $-m_3$ gives a surjective homomorphism
\begin{equation}\label{eq:H-to-Z2}
H\to\Z^2,\qquad p\mapsto(1,0),\quad q\mapsto(0,1),\qquad\text{with kernel }\FSym(Y).
\end{equation}

On the first ray we write $i$ for the point $(1,i)$ and $(i\ j)$ for the transposition of the points $i$ and $j$.

\begin{lemma}[Generators of $H$]\label{group:gen}
\leavevmode
\begin{enumerate}[label=(\alph*),itemsep=1pt]
\item For $k\ge0$ the element $\sigma_k=p^ksp^{-k}$ is the transposition $(k{+}1\ \,k{+}2)$; in particular $s=(1\ 2)$.
\item $H$ is generated by $p$ and $q$.
\end{enumerate}
\end{lemma}

\begin{proof}
(a) Composing from right to left, the letters $q^{-1},p^{-1},q,p$ of $s$ move $(1,1)$ through $(3,1),(3,1),(1,1)$ to $(1,2)$, and $(1,2)$
through $(1,1),(2,1),(2,1)$ to $(1,1)$. They move $(1,j)$ with $j\ge3$ through $(1,j-1),(1,j-2),(1,j-1)$ back to $(1,j)$, a point $(2,j)$
through $(2,j),(2,j+1),(2,j+1)$ back to $(2,j)$, and a point $(3,j)$ through $(3,j+1),(3,j+1),(3,j)$ back to $(3,j)$. So $s=(1\ 2)$. Since
$p^k$ maps the points $1,2$ to $k+1,k+2$, the conjugate $\sigma_k$ is $(k{+}1\ \,k{+}2)$.

(b) Let $H'=\langle p,q\rangle$. By (a), $H'$ contains the adjacent transpositions of ray~$1$, hence every finitary permutation of $Y$ supported
on ray~$1$. Let $D\subset Y$ be finite and choose $k,l$ larger than the heights of its points. The element $g=p^kq^l$ maps $D$ into ray~$1$:
$q^l$ maps ray~$1$ into itself and the points of height at most $l$ on ray~$3$ into ray~$1$, and fixes ray~$2$; then $p^k$ maps ray~$1$
into itself and the points of height at most $k$ on ray~$2$ into ray~$1$. So for a permutation $\pi$ supported in $D$, the conjugate
$g\pi g^{-1}$ is supported on ray~$1$ and lies in $H'$, and so does $\pi$. Hence $\FSym(Y)\le H'$, and since $H'$ maps onto $\Z^2$
under~\eqref{eq:H-to-Z2}, whose kernel is $\FSym(Y)$, we get $H'=H$.
\end{proof}

Johnson~\cite{Joh99} gave a finite presentation of $H$ on the generators $p$, $q$ and $\alpha=[p,q]$; see also~\cite[Theorem~2.8]{Lee12}.
Eliminating $\alpha$ gives the presentation
\begin{equation}\label{eq:PH}
P_H=\bigl\langle p,q\bigm|\ s^2,\ (sps\,p^{-1})^3,\ [s,p^2sp^{-2}],\ psp^{-1}qs^{-1}q^{-1}\bigr\rangle,\qquad s=[p,q].
\end{equation}
Its relators hold in $H$ by \cref{group:gen}(a): $s=(1\ 2)$ and $psp^{-1}=(2\ 3)$, so $sps\,p^{-1}$ is a $3$-cycle; $s$ commutes with
$p^2sp^{-2}=(3\ 4)$; and $qsq^{-1}=(2\ 3)=psp^{-1}$, because $q$ acts on ray~$1$ as $p$ does. We do not use Johnson's theorem:
\cref{app:dehn} reproves that $P_H$ presents $H$, with an explicit bound on its Dehn function (\cref{cor:dehn-qp}).

For a set $W$ of words, a word $w$ has \emph{area at most $k$ in $W$} if it is freely equal to a product of $k$ conjugates of elements of
$W^{\pm1}$; the \emph{area} $\Area_W(w)$ is the least such $k$, and $\infty$ if there is none. Let $\dH(n)$ be the largest area, with
respect to the four relators of $P_H$, of a word of length at most $n$ in $p^{\pm1},q^{\pm1}$ that represents $1$ in $H$. This is the
Dehn function of $P_H$; it is nondecreasing. Two bounds on it are available.
\begin{enumerate}[label=(H\arabic*),leftmargin=3.2em]
\item\label{H-exp} $\dH(n)\le2^{(\log n+6)^2}$ for all $n\ge1$. This is \cref{cor:dehn-qp}, proved in \cref{app:dehn} from Johnson's
presentation alone, by splitting in half Lee's recursion for commutators of far-apart transpositions~\cite[Section~4.2]{Lee12}. Lee's own
bound is exponential~\cite[Theorem~D]{Lee12}.
\item\label{H-poly} $\dH(n)\le C_Hn^{p_H}$ for an absolute constant $C_H$, where $p_H=4+\log107\approx10.74$. This is proved
in~\cite[Theorems~3.2 and~3.3]{DM-H3} for Johnson's presentation, through the bound $135M^{\log107}$ on the area of the commutator of two
transpositions at distance $M$, and it passes to $P_H$ by \cref{dehn:PJ}. (Permutations act on the right in~\cite{DM-H3}, and its
generators are the inverses, as maps, of ours; since $g\mapsto g^{-1}$ turns that product into composition, the two presentations match
letter by letter and have the same areas.)
\end{enumerate}
Every theorem of this paper uses only~\ref{H-exp}. The bound~\ref{H-poly} rests on an unrefereed preprint and appears only in remarks,
where it improves factors $2^{-O((\log\log(1/\eps))^2)}$ to powers of $\log(1/\eps)$; those remarks write $\kappa=2+p_H=6+\log107$.

\subsection{The affine group of configurations}\label{sec:group-affine}

Let $V=\F_2^{(Y)}$ be the space of finitely supported functions $Y\to\F_2$, the \emph{configurations}, with basis $(e_y)_{y\in Y}$. We
identify a configuration with its support, a finite set of points. The group $H$ acts faithfully on $V$ by permutation matrices,
$g(e_y)=e_{g(y)}$, and we identify $H$ with its image in $\GL(V)$. Let $\GL_{\rm fin}(Y)$ be the group of linear automorphisms of $V$ that
fix all but finitely many basis vectors. It is the directed union of the finite groups $\GL(\F_2^D)$ over finite $D\subset Y$, each extended
by the identity on the other basis vectors. It is normalized by $H$, and $H\cap\GL_{\rm fin}(Y)=\FSym(Y)$, because a permutation matrix
fixes all but finitely many basis vectors exactly when the permutation is finitary. Put
\[
L=\GL_{\rm fin}(Y)\,H\le\GL(V),\qquad A=V\rtimes L ,
\]
so that $A$ is the group of affine maps $v\mapsto g(v)+u$ of $V$ with $g\in L$ and $u\in V$. Write $\tau_u$ for the translation $v\mapsto
v+u$, and for distinct $y,y'\in Y$ let $E_{y\leftarrow y'}$ be the transvection that sends $e_{y'}$ to $e_{y'}+e_y$ and fixes the other
basis vectors. For every permutation $g$ of $Y$, acting on $V$ by permutation matrices, and every $M\in\GL(V)$,
\begin{equation}\label{eq:affine-id}
g\tau_{e_y}g^{-1}=\tau_{e_{g(y)}},\qquad gE_{y\leftarrow y'}g^{-1}=E_{g(y)\leftarrow g(y')},\qquad M\tau_uM^{-1}=\tau_{M(u)} .
\end{equation}
On the first ray we write $e_k=e_{(1,k)}$ and $\tau_k=\tau_{e_k}$, and put
\[
t=\tau_1,\qquad d=E_{(1,2)\leftarrow(1,1)} .
\]
Thus $t$ toggles the point $(1,1)$ in a configuration, and $d$ toggles $(1,2)$ in every configuration that contains $(1,1)$. Linear maps fix
the zero configuration $0$.

\begin{lemma}\label{lem:A-basic}
\leavevmode
\begin{enumerate}[label=(\alph*),itemsep=1pt]
\item $L$ is generated by $p,q,d$, and $A$ is generated by $p,q,t,d$.
\item $A$ acts doubly transitively on $V$, and the stabilizer of $0$ is $L$.
\end{enumerate}
\end{lemma}

\begin{proof}
(a) The subgroup $\langle p,q,d\rangle$ contains $H$ by \cref{group:gen}(b). Since $\FSym(Y)\le H$ acts doubly transitively on $Y$, the
second identity of~\eqref{eq:affine-id} puts every transvection $E_{y\leftarrow y'}$ into it. Transvections generate $\GL(\F_2^D)$ for every
finite $D$, by row reduction, so $\langle p,q,d\rangle$ contains $\GL_{\rm fin}(Y)$, hence $L$. By the first identity of~\eqref{eq:affine-id}
the $H$-conjugates of $t$ are all the $\tau_{e_y}$, which generate the translations; so $\langle p,q,t,d\rangle=A$.

(b) An affine map that fixes $0$ is linear, so the stabilizer of $0$ is $L$. Given $v\ne w$ in $V$, the translation $\tau_v$ carries $(v,w)$
to $(0,v+w)$ with $v+w\ne0$. Choose a finite $D\subset Y$ containing $(1,1)$ and the support of $v+w$; some $M\in\GL(\F_2^D)\le\GL_{\rm
fin}(Y)$ sends $v+w$ to $e_1$, so $M\tau_v$ carries $(v,w)$ to $(0,e_1)$. Hence $A$ is transitive on ordered pairs of distinct
configurations.
\end{proof}

\subsection{Pauli lamps}\label{sec:group-lamps}

Let $D_4=\langle x,z\mid x^2,z^2,(xz)^4\rangle$ be the dihedral group of order eight. Its centre is $\{1,(xz)^2\}$, and $(xz)^2=[x,z]$
because $x$ and $z$ are involutions. The assignment $x\mapsto\sigma_X$, $z\mapsto\sigma_Z$ to the Pauli matrices is a faithful
representation that sends $(xz)^2$ to $-I$; we call $D_4$ the \emph{Pauli lamp}.

For $a,b\in\F_2^{(V)}$, finitely supported functions $V\to\F_2$, put $\langle b,a\rangle=\sum_{v\in V}b(v)a(v)\in\F_2$. Let $\EV$ be the
set $\F_2\times\F_2^{(V)}\times\F_2^{(V)}$ with the product
\begin{equation}\label{eq:EV-product}
(c,a,b)\,(c',a',b')=\bigl(c+c'+\langle b,a'\rangle,\ a+a',\ b+b'\bigr).
\end{equation}
This is a group: by bilinearity both bracketings of a triple product have first coordinate
$c+c'+c''+\langle b,a'\rangle+\langle b,a''\rangle+\langle b',a''\rangle$, the identity is $(0,0,0)$, and the inverse of $(c,a,b)$ is
$(c+\langle b,a\rangle,a,b)$. Put
\[
J=(1,0,0),\qquad x_v=(0,\mathbf 1_v,0),\qquad z_v=(0,0,\mathbf 1_v)\qquad(v\in V),
\]
where $\mathbf 1_v$ is the indicator function of $\{v\}$, and for $a,b\in\F_2^{(V)}$ let $X^a=\prod_vx_v^{a(v)}$ and $Z^b=\prod_vz_v^{b(v)}$.
By~\eqref{eq:EV-product} the $x_v$ commute with each other, as do the $z_v$; $x_v$ commutes with $z_w$ for $v\ne w$; and
\[
z_vx_v=(1,\mathbf 1_v,\mathbf 1_v)=J\,x_vz_v,\qquad (x_vz_v)^2=J,\qquad J^cX^aZ^b=(c,a,b).
\]
So every element of $\EV$ has the unique \emph{normal form} $J^cX^aZ^b$, and~\eqref{eq:EV-product} is the rule
$Z^bX^{a'}=J^{\langle b,a'\rangle}X^{a'}Z^b$. The element $J$ is central in $\EV$, since $\langle b,0\rangle=\langle0,a\rangle=0$, and $J^2=1$.

\begin{lemma}\label{lem:EV-central-product}
The homomorphism $\bigoplus_{v\in V}D_4\to\EV$ that sends $x$ and $z$ in the summand at $v$ to $x_v$ and $z_v$ is surjective, and its
kernel is the subgroup $Z_{\rm ev}$ of products of an even number of the central elements $J_v=(xz)^2_v$. Thus $\EV$ is the central product
of Pauli lamps, one at each configuration, with all centres identified to $J$.
\end{lemma}

\begin{proof}
The relations of $D_4$ hold for $x_v,z_v$, since $(x_vz_v)^4=J^2=1$, and elements attached to distinct sites commute, so the homomorphism
exists. Its image contains all $x_v$, $z_v$ and $J=(x_0z_0)^2$, so it is onto by the normal form. Each $J_v$ maps to $J$, so $Z_{\rm ev}$
lies in the kernel. For a finite nonempty set $F\subset V$, the subgroup $\bigoplus_{v\in F}D_4$ has order $8^{|F|}$ and meets $Z_{\rm ev}$
in a group of order $2^{|F|-1}$, while its image $\{(c,a,b):a,b\text{ supported in }F\}$ has order $2\cdot4^{|F|}=8^{|F|}/2^{|F|-1}$. So the
kernel on $\bigoplus_{v\in F}D_4$ is exactly $Z_{\rm ev}\cap\bigoplus_{v\in F}D_4$, and since $F$ is arbitrary the kernel is $Z_{\rm ev}$.
\end{proof}

The group $A$ permutes $V$, and it acts on $\EV$ by moving sites and fixing $J$:
\[
g\cdot(c,a,b)=\bigl(c,\ a\circ g^{-1},\ b\circ g^{-1}\bigr),\qquad\text{so that}\qquad g\cdot x_v=x_{g(v)},\quad g\cdot z_v=z_{g(v)},\quad
g\cdot J=J .
\]
This is an action by automorphisms, because $g$ is a bijection of $V$ and so $\langle b\circ g^{-1},a'\circ g^{-1}\rangle=\langle
b,a'\rangle$. Under \cref{lem:EV-central-product} it is induced by the permutation of the summands of $\bigoplus_VD_4$, which preserves
$Z_{\rm ev}$.

\subsection{The groups \texorpdfstring{$\GP$}{G\_P} and \texorpdfstring{$\GPp$}{G\_P+}}\label{sec:group-GP}

\begin{definition}[The Pauli configuration-lamp group]\label{def:GP}
The \emph{Pauli configuration-lamp group} is $\GP=\EV\rtimes A$. We write $x=x_0$ and $z=z_0$ for the Pauli pair at the zero
configuration, so that $J=(xz)^2=[x,z]$, and $p,q,t,d\in A$ are as above.
\end{definition}

In $\GP$ we have $gx_vg^{-1}=x_{g(v)}$, $gz_vg^{-1}=z_{g(v)}$ and $gJg^{-1}=J$ for $g\in A$. Equivalently, $\GP$ is the quotient of the
permutational wreath product $D_4\wr_VA=\bigl(\bigoplus_VD_4\bigr)\rtimes A$ by the central subgroup $Z_{\rm ev}$, which is normal
because $A$ preserves it.

\begin{proposition}\label{prop:J-central}
$J$ is a central involution of $\GP$ and $J\ne1$. It commutes with $p,q,t,d,x,z$, and these six elements generate $\GP$.
\end{proposition}

\begin{proof}
$J$ is central in $\EV$ and fixed by the action of $A$, so it is central in $\EV\rtimes A$; $J^2=1$ by~\eqref{eq:EV-product}, and
$J=(1,0,0)\ne(0,0,0)$. By \cref{lem:A-basic}(a) the subgroup generated by the six elements contains $A$. For $v\in V$ the translation
$\tau_v$ sends $0$ to $v$, so the subgroup contains $x_v=\tau_vx\tau_v^{-1}$, $z_v=\tau_vz\tau_v^{-1}$ and $J=(xz)^2$, which generate $\EV$
by the normal form.
\end{proof}

\begin{remark}[Finite presentation]\label{rem:GP-fp}
The group $\GP$ is finitely presented. By~\cite[Proposition~7.1]{DM-lamps} the group $A$ is finitely presented, so by
\cref{lem:A-basic}(a) it has a finite presentation $\mathcal R_A$ on $p,q,t,d$. Adding the generators $x,z$ and the relators $x^2$, $z^2$,
$(xz)^4$, $[\ell,u]$ for $\ell\in\{p,q,d\}$ and $u\in\{x,z\}$, $[u,tu't^{-1}]$ for $u,u'\in\{x,z\}$, and $[(xz)^2,t]$, which are the
lamp words of the certificate $\RP$ of \cref{def:RP}, gives a presentation of $\GP$. Indeed, the stabilizer $L=\langle p,q,d\rangle$ of $0$
centralizes $\langle x,z\rangle$, so the conjugate $g\langle x,z\rangle g^{-1}$ depends only on $g(0)$; double transitivity
(\cref{lem:A-basic}(b)) carries the commutation of the pairs at $0$ and $e_1$, and the identification of their centres, to every two
distinct sites; and \cref{lem:EV-central-product} identifies the group so presented with $\GP$. This paper never uses finite
presentability. All profiles are taken with respect to explicit finite sets of words valid in $\GP$ that contain $\RP$, and Theorem~B holds
for every such set, hence in particular for every finite presentation of $\GP$ on $p,q,t,d,x,z$.
\end{remark}

\begin{proposition}[Amenability]\label{prop:GP-amenable}
Let $\pi:\GP\to\Z^2$ be the composite $\GP\to A\to L\to L/\GL_{\rm fin}(Y)\cong H/\FSym(Y)\cong\Z^2$, the last map being~\eqref{eq:H-to-Z2}.
Its kernel $\EV\rtimes\bigl(V\rtimes\GL_{\rm fin}(Y)\bigr)$ is locally finite. Hence $\GP$ is locally finite-by-$\Z^2$; in particular it is
elementary amenable, amenable, sofic and hyperlinear.
\end{proposition}

\begin{proof}
The map $A\to L$ is the linear part, with kernel $V$. Since $L=\GL_{\rm fin}(Y)H$ with $\GL_{\rm fin}(Y)$ normal, the inclusion induces an
isomorphism $H/(H\cap\GL_{\rm fin}(Y))\cong L/\GL_{\rm fin}(Y)$, and $H\cap\GL_{\rm fin}(Y)=\FSym(Y)$; so $\pi$ is well defined and onto, and
its kernel is as stated. The group $\EV$ is locally finite: a finitely generated subgroup lies in the subgroup of elements $(c,a,b)$ with
$a,b$ supported on a fixed finite set $F\subset V$, which has order $2\cdot4^{|F|}$. The group $V\rtimes\GL_{\rm fin}(Y)$ is the directed
union of the finite groups $\F_2^D\rtimes\GL(\F_2^D)$ over finite $D\subset Y$, so it is locally finite. An extension of a locally finite
group $K$ by a locally finite group is locally finite: a finitely generated subgroup $\Lambda$ has finite image modulo $K$, so $\Lambda\cap K$
has finite index in $\Lambda$, is therefore finitely generated, and is finite. So $\ker\pi$ is locally finite.

The class of elementary amenable groups is the smallest class containing the finite and the abelian groups and closed under subgroups,
quotients, extensions and directed unions~\cite{Cho80}. A locally finite group is the directed union of its finite subgroups, so $\ker\pi$
is elementary amenable, and so is the extension $\GP$ of $\ker\pi$ by $\Z^2$. Elementary amenable groups are amenable, amenable groups are
sofic~\cite{Gro99,Wei00}, and sofic groups are hyperlinear~\cite{ES05}; see~\cite{Pes08}.
\end{proof}

\paragraph{The ray swap.} Let $\iota$ be the permutation of $Y$ that exchanges rays $2$ and $3$: $\iota(1,j)=(1,j)$, $\iota(2,j)=(3,j)$ and
$\iota(3,j)=(2,j)$. It is not in $H$, but it normalizes $H$ (conjugation by $\iota$ exchanges the translation lengths on rays $2$ and $3$)
and $\FSym(Y)$. As a permutation matrix on $V$ it normalizes $\GL_{\rm fin}(Y)$, hence $L$, and by~\eqref{eq:affine-id} it normalizes the
translations, hence $A$. It does not lie in $A$: it fixes $0$, so it would lie in $L$ (\cref{lem:A-basic}(b)), and if $\iota=Mh$
with $M\in\GL_{\rm fin}(Y)$ and $h\in H$, then $M=\iota h^{-1}$ is a permutation matrix that fixes all but finitely many basis vectors, so
$\iota\in\FSym(Y)H=H$. So $A^+=\langle A,\iota\rangle=A\rtimes\langle\iota\rangle$ is a group of affine maps of $V$, with
$\langle\iota\rangle\cong\Z/2$; it acts on $\EV$ by moving sites, by the formula above, and $\iota$ fixes $0$. Put
\[
\GPp=\EV\rtimes A^+=\GP\rtimes\langle\iota\rangle .
\]

\begin{lemma}\label{lem:GPp}
In $\GPp$, $\iota p\iota=q$, and $\iota$ commutes with $t$, $d$, $x$, $z$ and $J$. So $\GPp$ is generated by $p,\iota,t,d,x,z$, the element
$J$ is central in $\GPp$, and $\GPp$ is elementary amenable, hence hyperlinear.
\end{lemma}

\begin{proof}
Since $\iota=\iota^{-1}$, the element $\iota p\iota$ is the conjugate $\iota p\iota^{-1}$. We compare it with $q$ as permutations of $Y$,
composing from right to left. On ray~$1$, $\iota p\iota(1,j)=\iota(1,j+1)=(1,j+1)$. On ray~$3$, $\iota p\iota(3,1)=\iota p(2,1)=(1,1)$ and
$\iota p\iota(3,j)=\iota(2,j-1)=(3,j-1)$ for $j\ge2$. On ray~$2$, $\iota p\iota(2,j)=\iota p(3,j)=\iota(3,j)=(2,j)$. This is the
definition of $q$. Since $\iota$ fixes the points $(1,1)$ and $(1,2)$, the identities~\eqref{eq:affine-id} give $\iota t\iota=t$ and
$\iota d\iota=d$. Since $\iota$ is linear it fixes the site $0$, so it commutes with $x=x_0$ and $z=z_0$, and it fixes $J$. As $J$ is
central in $\GP$ (\cref{prop:J-central}) and commutes with $\iota$, it is central in $\GPp$. By \cref{prop:J-central} and $q=\iota p\iota$,
the elements $p,\iota,t,d,x,z$ generate $\GPp$. Finally $\GPp$ is an extension of $\GP$ by $\Z/2$, so it is elementary amenable, hence
hyperlinear, as in \cref{prop:GP-amenable}.
\end{proof}

Finite-dimensional models of $\GPp$ in which $J$ acts as $-I$ exactly are constructed explicitly in \cref{sec:upper}.
